\documentclass[10pt,a4paper]{article}
\usepackage[T1]{fontenc}
\usepackage[utf8]{inputenc}
\usepackage{lmodern,microtype}
\usepackage[margin=18mm,headheight=15pt,footskip=10mm]{geometry}
\usepackage{tikz}
\usetikzlibrary{arrows.meta,positioning}
\usepackage{xcolor,enumitem,fancyhdr,hyperref}
\definecolor{ink}{HTML}{173347}
\definecolor{soft}{HTML}{F2F5F7}
\definecolor{edge}{HTML}{8496A3}
\hypersetup{colorlinks=true,urlcolor=ink,linkcolor=ink,pdftitle={HOF analysis - co-generation plant},pdfauthor={AI-assisted draft for Arici, Ferraboli and Brusamarello}}
\pagestyle{fancy}\fancyhf{}
\fancyhead[L]{\small\textcolor{ink}{Technology Risk Governance | Minor HOF}}
\fancyhead[R]{\small Independent Codex draft | 25 Nov 2026}
\fancyfoot[L]{\scriptsize F: case fact \quad T: taught interpretation \quad I: proposed inference}
\fancyfoot[R]{\small\thepage/8}
\renewcommand{\headrulewidth}{0.3pt}
\setlength{\parindent}{0pt}\setlength{\parskip}{5pt}
\setlist[itemize]{leftmargin=14pt,itemsep=3pt,topsep=3pt}
\newcommand{\evtitle}[3]{\vspace*{-3mm}{\LARGE\bfseries\color{ink}Event #1: #2}\par{\large #3}\par\vspace{2mm}}
\newcommand{\sect}[1]{\par\vspace{2mm}{\bfseries\color{ink}#1}\par}
\newcommand{\tagline}[1]{{\small\textcolor{ink}{#1}}\par}
\newcommand{\sourcebar}[1]{\vfill{\footnotesize\textbf{Source key.} #1}}
\newif\ifhumanpath\humanpathtrue
\newcommand{\reason}[8]{%
\begin{center}\begin{tikzpicture}[x=1cm,y=1cm,>=Stealth,
 box/.style={draw=edge,rounded corners=1mm,fill=soft,align=left,text width=4.8cm,inner sep=7pt,font=\small},
 top/.style={box,text width=8cm},
 event/.style={box,fill=ink,text=white,text width=4.6cm,align=center,font=\small\bfseries},
 conn/.style={->,thick,draw=ink},context/.style={->,dashed,draw=edge}]
\node[top] (org) at (0,0) {\textbf{Organisational vulnerabilities / cross-path latent conditions}\\#1};
\node[font=\bfseries\small,text=ink] at (-5,-2.1) {Human error path};
\node[font=\bfseries\small,text=ink] at (5,-2.1) {System failure path};
\node[box] (local) at (-4.9,-4.2) {\textbf{Local working conditions}\\#2};
\node[box] (latent) at (4.9,-4.2) {\textbf{Latent errors}\\#3};
\node[box] (active) at (-4.9,-8.4) {\textbf{Active errors}\\#4};
\node[box] (fail) at (4.9,-8.4) {\textbf{Failures of defences, barriers and safeguards}\\#5};
\node[event] (ev) at (0,-13.1) {#6};
\node[box,text width=4.1cm] (human) at (-5.7,-13.1) {\textbf{Personal / external factors}\\#7};
\node[box,text width=4.1cm] (system) at (5.7,-13.1) {\textbf{Local / external factors}\\#8};
\draw[conn] (org.south) -- (0,-2.1) -- (-2,-2.1) -- (local.north east);
\draw[conn] (org.south) -- (0,-2.1) -- (2,-2.1) -- (latent.north west);
\draw[conn] (local) -- (active);\draw[conn] (latent) -- (fail);
\ifhumanpath\draw[conn] (active.south east) -- (ev.north west);\fi
\draw[conn] (fail.south west) -- (ev.north east);
\draw[context] (human.east) -- (ev.west);
\draw[context] (system.west) -- (ev.east);
\end{tikzpicture}\end{center}}
\begin{document}
\evtitle{1}{Start of transient depressurisation}{02:00 | Reason template | Required classroom event}
\tagline{Group: Michele Attilio Arici; Andrea Ferraboli; Michelangelo Brusamarello.}
\tagline{Working status: AI-assisted draft for substantive group revision; additional events and interpretations remain proposals.}
\reason
{Control-system design errors and electrical-circuit design/maintenance weaknesses [T, S10]. They undermine both reliable information and operator diagnosis [I].}
{Poor automatic control/HMI; no emergency procedures [T, S10; F, C\,\S3.3--3.4].}
{Pressure-transducer location exposed to freezing; insufficient control-logic redundancy [T, S10]. No abnormal-measure warning [F, C\,\S3.3].}
{Operators did not understand the DCS automatic routine: \textit{mistake} [T, S10]. This explains failure to recognise the developing state, rather than the initiating sensor fault.}
{Pressure transducer fails and supplies a false high-pressure signal; automatic controls act on it [F, C\,\S2; T, S10].}
{EVENT 1\\False 1.76 MPa indication triggers depressurisation although actual pressure was normal [F].}
{Night shift [F]. Low production and low attention are classroom interpretations [T, S10], not established individual behaviour.}
{Low winter temperature; freezing of signal piping [F]. Temperature of -5\,$^{\circ}$C is from the teaching solution [T, S10].}
\tagline{Read causal arrows as contributory mechanisms, not quantified probabilities. Dashed arrows show contextual influences.}
\sourcebar{C: Trucco, \textit{Accident in a co-generation plant}, 2014, pp.6--10. S: Class 8 presentation, PDF pp.7--10. R: Class 7 HOF, PDF pp.9,14,19.}
\newpage
\evtitle{1}{Justification and scope}{A failed information barrier, compounded by loss of situational awareness}
\sect{Evidence and causal allocation}
\textbf{F.} C\,\S2 (pp.6--8) places the transducer outdoors, describes freezing during an unusually cold but \emph{not exceptional} night, and records the erroneous 1.76 MPa indication against a normal system pressure near 0.95 MPa. The control system reduced pressurising-pump speed, opened relief/expansion-tank valves and stopped main feed pumps. The equipment responded to the wrong information; the case does not establish that all these actuators themselves malfunctioned.

\textbf{T.} S10 separates the failed transducer (failed defence) from its exposed location and missing control redundancies (latent errors). The top box of the teaching solution is labelled \emph{latent error}, whereas the generic Reason template labels it \emph{organisational vulnerabilities}. We retain its cross-path function and make the terminology explicit. This follows the lecture's allocation by \emph{scope of influence}: cross-path conditions go at the top; conditions specific to one path remain on that path (L, discussion 42:25--43:35 / 00:00--01:45).

\textbf{F/T.} C\,\S3.3--3.4 (p.9) establishes no abnormal-measure alarm, a poorly designed heater-supply monitoring system, layered HMI windows without overview and no emergency procedures. S10 calls the human active error a \emph{mistake}: operators misunderstood the DCS routine. The subsequent restart of circulation pumps supports that classroom interpretation (L, 18:19--24:15); it is not an independent observation of their mental state.
\sect{Why these are distinct categories}
The transducer defect initiates the technical sequence; misunderstanding compromises the opportunity for recovery. A location/control-design weakness predates the incident, so it is latent. The HMI and unavailable procedures are the conditions experienced by operators. The inferred failure of understanding concerns intention/diagnosis, so \emph{mistake} is consistent with Reason's taxonomy (R14), rather than an execution slip. No claim is made that an operator should have prevented an instantaneous automatic response.
\sect{Uncertainties and proportionate implications}
Low attention is retained solely as a taught interpretation. Night shift alone does not demonstrate fatigue, sleep or misconduct; such claims are excluded. C\,\S3.6 and \S3.8 explicitly lack detailed shift/team information. The awkward wording about thermo-hydraulic knowledge in C\,\S3.1 is interpreted using the classroom explanation and C\,\S3.7, which directly states inadequate emergency skills.

\textbf{I, proposed improvements.} Protect and supervise the sensor/heating circuit; validate critical readings using suitably independent information; expose control mode, pressure plausibility and actuator states together; test emergency procedures through realistic drills. Independence/common-cause checks are needed before adding sensors. These address the identified barriers and do not reduce the lesson to operator blame.
\sourcebar{L: supplied \texttt{trascrizione\_lezione\_caso.txt}, sections on Event 1 and methodological questions. L is a structured guide derived from a reordered video, not a verbatim transcript; C and S take priority.}
\newpage
\evtitle{3}{Flooding of the pumping room}{02:41 | Reason template | Required classroom event}
\reason
{Alarm-system design and overflow-emergency logic [T, S11]. Inadequate alignment of containment capacity with detection and intervention [I, following L].}
{High-level alarm buried among lower-priority alarms; layered displays obscure overall plant condition [F, C\,\S3.3; T, S11].}
{Discharge line/tank undersized; insufficient isolation of the pump room [T, S11]. Inadequate emergency discharge capacity is explicit in the case [F].}
{Omitted valve closure: \textit{mistake}; high water level not noticed: \textit{slip} [T, S11]. Late closure is documented separately as Event 4 [F].}
{Storage/overflow system cannot contain or discharge the incoming water; boiling water reaches the pumping room [F, C\,\S2; T, S11].}
{EVENT 3\\Storage-tank overflow and flooding of the pumping room [F].}
{Limited emergency response time [F, C\,\S3.5]. No event-specific personal state is established; no individual misconduct inferred.}
{Earlier manual activation of circulation pumps (Event 2) [F; external input in the event-specific teaching solution, T, S11].}
\tagline{Event 2 is a prior human action in the accident sequence, but an external input to this Event 3 boundary. The classification changes with the event boundary.}
\sourcebar{C, \S2 pp.7--8; \S3.3--3.5 pp.9--10. S11, official classroom analysis of Event 3. Generic Reason scheme: S7 / R19.}
\newpage
\evtitle{3}{Justification and scope}{Containment capacity must match the time needed for an effective response}
\sect{Evidence and causal allocation}
\textbf{F.} C\,\S2, Event 3 (p.8), states that maximum storage-tank level led to discharge and that the emergency system could not fully handle the flow; boiling water flooded the pumping room. C\,\S3.3 (p.9) states that a high-level warning existed but was dispersed among less important alarms and that layered HMI windows prevented an overview. This is a failure of effective warning presentation, not evidence that no level alarm was installed.

\textbf{T.} S11 assigns deficient pump-room isolation and discharge-line/tank sizing to latent technical errors, and overflow to failure in defence. It identifies omitted valve closure as a \emph{mistake}, and failure to notice high water level as a \emph{slip}. We reproduce those taught classifications. The case directly documents late closure in Event 4, but does not independently establish whether every operator formed an incorrect intention or missed an execution cue. The mental-error labels therefore remain T, not F.

\textbf{T/I.} L (29:00--29:49 and 00:50--01:45) emphasises alignment of technical sizing with organisational response. Finite capacity is useful only when detection, diagnosis and closure can occur before it is exhausted. Placing the joint design weakness at the top connects both causal paths without claiming a documented management decision or a numerical capacity deficit.
\sect{Boundary and taxonomy}
The manual pump restart is external to \emph{this sub-event}, as in S11. Its consequences increased the demand on discharge/containment; the full sequence still includes the human decision in Event 2. Generic training shortcomings exist in C\,\S3.7, but are not inserted as the explanation for every box. L's SRK discussion (35:05--39:50) distinguishes basic monitoring/closure from the deeper diagnosis required in Event 1, and allows justified alternative allocations. The present diagram stays close to S11.
\sect{Uncertainties and proportionate implications}
No tank volume, discharge flow, time-to-overflow or individual alarm-response record is supplied. Classroom claims about flow multiples are therefore not used as measured design data. The report does not infer staffing numbers, fatigue or a particular maintenance failure.

\textbf{I, proposed improvements.} Assess credible inflows jointly with operator response time; segregate drainage from vulnerable equipment; prioritise high-level alarms and provide an integrated view; define a safe valve/pump sequence and verify it in drills. Improvements must be tested together: a larger tank alone cannot guarantee safe recovery, and a procedure alone cannot compensate for an ineffective warning or insufficient physical capacity.
\sourcebar{C pp.8--10; S11; R14--19. L Event 3 and methodological questions, cited by supplied guide section/minute range. The case's repeated \S3.4 numbering is preserved when locating sources.}
\newpage
\evtitle{5}{Loss of power and remote control}{03:20 | Reason template | Proposed additional event A}
\humanpathfalse
\reason
{Integrated protection design did not preserve an effective monitoring/control layer under flooding and steam conditions [I, based on F]. No emergency procedures; inadequate emergency readiness [F, C\,\S3.4,3.7].}
{Remote control becomes unavailable; recovery must use local/manual arrangements [F, C\,\S1.2--1.3,2]. Emergency staffing is inadequate [F, C\,\S3.1].}
{Fire protection cannot distinguish vapour from smoke [F]. Protective-trip and control-power dependencies allow a common loss of operational capability [I].}
{No distinct operator error is established for triggering Event 5 [F boundary]. Human recovery difficulties are consequences/conditions, not an invented unsafe act.}
{Fire protection trips on steam, cuts electrical power, stops pumps and removes remote control [F, C\,\S2]. Expected continuity of essential control is not achieved [I].}
{EVENT 5\\Fire-protection intervention produces complete power cut-off and remote-control outage [F].}
{Night shift and limited emergency time [F]. No documented individual attention or physiological state.}
{Boiling-water flooding and vapour-saturated atmosphere from earlier events [F]. No actual fire is reported in this sequence.}
\humanpathtrue
\tagline{Human branch: recovery context only; no established human-error link to the trip. Selection proposal: complements Events 1 and 3 through safeguard interaction; not yet confirmed by the group.}
\sourcebar{C, \S1.2--1.3 pp.5--6; \S2 pp.6--8; context pp.9--10. R5--9,19: layered defences and active/latent conditions. Event 5 classification is our proposed analysis, not a classroom solution.}
\newpage
\evtitle{5}{Justification and scope}{A protective action creates a new loss of control capability}
\sect{Evidence and causal allocation}
\textbf{F.} C\,\S2 (p.8) directly links flooding and vapour saturation to fire-protection intervention. That system could not distinguish vapour from smoke and cut off electrical power; pumps stopped and remote control went out of service. These are the documented proximal mechanisms. The water/steam condition is an input from the preceding sequence, not a spontaneous fire or a cyber incident.

\textbf{F/I.} C\,\S2 (pp.6--7) also describes a backup supply intended to retain the minimum power for automatic-control settings during a fire trip. Its coexistence with the recorded remote-control outage reveals a continuity question. It does \emph{not} prove that a particular UPS, battery, wiring circuit or component failed. Retaining settings and sustaining remote operation may require different capabilities. The diagram therefore records the observed loss and a proposed dependency concern, without inventing its electrical failure mechanism.

\textbf{F/I.} C\,\S1.2 (p.5) defines manual operation as an emergency backup, rarely used because of complexity/difficulty. C\,\S1.3 (pp.5--6) describes separated control areas and local panels. Their contribution here is to recovery conditions after remote control is lost. Inadequate staffing, absence of emergency procedures and inadequate emergency skills are case facts (pp.9--10); their placement as cross-path organisational vulnerabilities is our causal interpretation.
\sect{Why no forced human-error entry}
The documented action initiating this sub-event is automatic. There is no evidence that an operator caused the detector to confuse steam and smoke or intentionally removed control power. Reason's two paths organise the evidence; they do not require an unsupported human mistake in every event. The human branch records constrained recovery and explicitly leaves the active-error conclusion unestablished. It should not be read as a separate proven cause of the trip.
\sect{Selection rationale and improvement implications}
\textbf{I.} Event 5 is proposed because it tests \emph{defence-in-depth} (R5--9): one protective layer should not inadvertently eliminate the next layer's usable capability under the same credible disturbance. It adds a barrier-interaction perspective beyond the initiating false measurement and overflow.

\textbf{I, proposed improvements.} Review steam/fire discrimination against actual fire-safety requirements; isolate essential control power and communications from flood/trip dependencies; test what backup supply can sustain; provide a safe, rehearsed local recovery path. Changes must preserve effective fire protection. A proposal to suppress fire trips or to keep all equipment energised in a genuine fire would not follow from this analysis. Specific hardware redundancy or battery autonomy cannot be prescribed without missing design information.
\sourcebar{Primary evidence: C pp.5--10, especially \S2 Event 5 on p.8. Method: R5--9,19. This additional-event analysis is I and requires independent group evaluation.}
\newpage
\evtitle{7}{Pump restart with protection bypassed}{04:39 | Reason template | Proposed additional event B}
\reason
{Priority given to heating continuity, conflicting restoration/safe-state goals, no emergency procedures and inadequate emergency competence [F, C\,\S3]. Their joint influence is inferred [I].}
{Remote control already lost; complex manual recovery; limited time; supervisor moving between separated rooms [F, C\,\S1.2,2,3.2,3.5].}
{Adequacy of bypass controls, restart permissions and safe-state verification is unestablished. An inadequate restart/bypass safety design is a hypothesis to investigate [I].}
{Manual restart while bypassing low-pressure protection, after three attempts [F]. Proposed diagnosis/planning \textit{mistake} [I]; a procedural violation is not established.}
{Low-pressure protection no longer inhibits the manually initiated restart [F]. Safe recovery is not ensured: further pressure reduction then violent pressurisation [F].}
{EVENT 7\\Recirculation pumps manually activated with low-pressure protection bypassed [F].}
{Night shift; mainly on-the-job training without periodic reinforcement; inadequate emergency skills [F]. No individual intent or fatigue evidence.}
{System already at low pressure after earlier transients; flooding, steam and previous power/control loss constrain recovery [F].}
\tagline{Selection proposal: examines intentional intervention under conflicting objectives. The final shock wave is Event 9, after further events; it is not equated with the restart itself.}
\sourcebar{C, Fig.3 p.7; \S2 Event 7 p.8; context pp.9--10. R13--16,19: violation, mistake and conditions. No additional-event solution from an excluded source was consulted.}
\newpage
\evtitle{7}{Justification and scope}{A bypass is documented; the operator's intent and the bypass design are not}
\sect{Evidence, interpretation and alternatives}
\textbf{F.} C\,\S2, Event 7 (p.8), states that manual pump activation bypassed low-pressure protection, that three attempts preceded activation, and that further pressure reduction was followed by violent pressurisation. Fig.3 (p.7) places this at 04:39. The bypass defeats a specific inhibition; it does not establish the complete control/protection architecture or prove that bypass hardware was defective.

\textbf{I.} A diagnosis/planning \emph{mistake} is a plausible classification: operators sought restoration without an adequate appreciation of the unsafe hydraulic state. Support comes from documented competing restoration/safety goals, emergency knowledge/experience limitations and absence of emergency procedures (C pp.9--10). Neither intention nor exact reasoning was recorded, so this remains a proposal. \emph{Violation} requires a known rule/practice deliberately departed from (R13); the deliberate bypass alone cannot establish awareness of a prohibited action. A slip/lapse is also not supported by a documented intended safe plan or execution record.

\textbf{I.} The technical latent-error box remains an investigation question. The record supports a bypassable protection, but not a conclusion that a second interlock was absent, that authorisation rules were ignored, or that design standards were breached. This restraint distinguishes the direct defence defeat from its unverified design causes. The upper cross-path conditions, by contrast, are grounded in the supplied organisational context.
\sect{Chronology and selection}
\textbf{F.} Timeline (C Fig.3): 02:00 E1 false-pressure/depressurisation; 02:04 E2 pump restart; 02:41 E3 overflow; 03:00 E4 late valve closure; 03:20 E5 fire trip/control loss; 03:45 E6 outlet-temperature rise; 04:39 E7 protection-bypassed restart; 05:15 E8 heat recovery started; 06:15 E9 shock wave. C p.8 specifies 90\,mm displacement and 3--4\,MPa peak. A conflicting classroom-guide phrase suggesting nearly one metre is rejected in favour of the case.

\textbf{I.} Event 7 is proposed alongside Event 5 to compare an automatic safeguard interaction with a human intervention. Intervening Event 8 and the evolving two-phase state preclude attributing Event 9 exclusively to one operator action. Proposed improvements are a verified safe-state/restart checklist, controlled and logged bypass authority, clear joint command across control rooms and hydraulic-emergency drills. Feasibility and any need for technical restart constraints require further design evidence.
\sect{Use status, sources and GenAI disclosure}
Prepared with Codex using delegated analysis of allowed teaching sources. AI selected the \emph{proposed} additional events, drafted classifications, reasoning and text, and generated the editable source/PDF; no human validation is claimed. The group must independently select/rework these judgements and review the text before submission. The 20 October lesson (line 221) expressly extends the GenAI principles to major and minor. Their restrictions require independent, substantive student work; a signature alone does not satisfy them.

\textbf{References.} C: P. Trucco (2014), \emph{Accident in a co-generation plant: analysis of the HOF involved}, pp.1--12. R: \emph{07\_SLIDES\_TRG\_HOF\_Class 7}, pp.5--19. S: \emph{08\_SLIDES\_TRG\_Class 8\_HOF Analysis case\_Presentation}, pp.7--12. L: supplied derived lesson guide, \texttt{trascrizione\_lezione\_caso.txt}, especially 13:53--15:40 (four events; one template page plus at most one justification page each). Original case/slides prevail over annotations and guide; detailed source paths and decisions accompany this report.
\end{document}
